
\setcounter{equation}{0}
\setcounter{Q}{0}


\begin{center}
\begin{minipage}{0.78\textwidth}
\includegraphics[width=\textwidth]{/home/jeff/Documents/IUTOrsay/1-Cours/4p13-Alliages/Cours-TD/Exos/Binaires-Lecture/CompDefini-Ex01/figures/Fig01}
\end{minipage}\hfill\begin{minipage}{0.18\textwidth}
\captionof{figure}{Diagramme binaire isobare, à une pression de \SI{1}{\bar}, du mélange \ce{Ce}-\ce{Al}.}
[La composition est donnée par la fraction molaire en cérium dans le mélange]
\end{minipage}
\end{center}

La figure ci-contre représente une partie du diagramme de phases solide-liquide du systèm binaire aluminium-cérium.

\Q Trois composés définis peuvent être repérés sur ce diagramme. Ecrire leurs formules.

\Q Indiqur la composition des domaines \textbf{I} à \textbf{XI} de ce diagramme. Que représente le point \textbf{J}.

\Q Donner l'allure de la courbe de refroidissement (température en fonction du temps) à partir de \SI{1000}{\celsius} d'un mélange de cérium et d'aluminium dont la fraction molaire de cérium est $0.75$. Pour chaque portion de courbe, préciser la nature de la ou des phases constituants le système.

Même question dans le cas d'un mélange cérium-aluminium ayant la composition du point \textbf{J}.

\Q Quelle masse d'un composé défini (que l'on précisera), peut-on obtenir pur, en partant du mélange du domaine \textbf{VII} de fraction molaire $x = 0.8$. \`{A} quel moment du refroidissement fait-il s'arrêter pour obtenir un composé pur ?

\paragraph*{Données}
\begin{align*}
M(\ce{Al}) &= 27 ~ \si{\gram\per\mole} ; \\
M(\ce{Ce}) &= 140 ~ \si{\gram\per\mole} .
\end{align*}